\section{Conclusiones}
\label{iv:conclusiones}
\begingroup
\setlength{\parskip}{0pt}
\setlength{\abovedisplayskip}{5pt}
\setlength{\belowdisplayskip}{5pt}

La estructura discreta del continuo desarrollada en este artículo
sostiene una realización excepcional que alcanza \(V^\natural\) y
una realización operatoria que conserva la dualidad compacta.
El vínculo está construido: el estado APP--TRIT--TPK transmite
orientación, hojas, residuos, cocientes, frontera y memoria a sus
coordenadas numéricas, al registro de incidencia y a las representaciones
de fase y acción. La contribución central consiste en relacionar estas
realizaciones mediante códigos, pegados, productos de vértice,
proyectores y operadores de entrelazamiento con dominios determinados.

La construcción conjunta del continuo proporciona el antecedente de
esas aplicaciones. Las mismas prolongaciones producen la medida de
historias, la identidad de Gram por hojas con su residuo terminal,
el defecto de cancelación, la incidencia rectangular y el transporte
natural del complejo de memoria. El ensamblaje conserva todas las
continuaciones y satisface \(D^2=0\) sobre sus generadores; su línea
determinante conserva grados y bases. Los pesos compatibles, la
isometría del transporte, el residuo de supervivencia y las condiciones
de una lectura determinante escalar mantienen la función concreta que
les asignan las pruebas. Estas operaciones correlacionadas constituyen
la estructura común anterior a las representaciones numéricas,
excepcionales y dimensionales.

% BEGIN R27_FRAMING_CONCLUSION_ES
La articulación aritmética e incidencial queda precisada por resultados
con dominios explícitos. Las leyes del defecto distinguen el refinamiento
del módulo, la repetición periódica y el aumento de dimensión. El operador
agregado conserva una subretícula integral de índice dos. En su bloque
cuadrático \(3H\), la unidad normalizada \(H\) tiene norma uno; la
recurrencia de Pell determina \(H^n\) y el bloque original evoluciona
como \(3^nH^n\).
La biyección entre pares de trits y dígitos nonádicos conserva el acarreo,
mientras la proyección del cubo ambiente tiene nueve elementos por fibra.
El alfabeto de \(42\) bloques, los empalmes regionales y las correlaciones
de registros ordenados hacen visible la información anterior a las
representaciones lineales.

% BEGIN R28_FRAMING_CONCLUSION_ES
Los cinco cierres regionales alcanzan el mínimo ocho en el problema
finito de cobertura de las cuatro aristas prescritas. El grafo conserva
las etiquetas de las emisiones y el cociente conserva sus incidencias
entre clases de fase. La inversión orientada hace impar al lector
\(\nu_{120}\) y mantiene par a \(\nu_{270}\); estas identidades
explicitan la respuesta de la memoria al cambio de orientación.
% END R28_FRAMING_CONCLUSION_ES

% BEGIN S27_FRAMING
La comparación regional determina tres configuraciones polarizadas:
las estrellas de \(B_K\), \(B_\varphi\) y \(B_{\rm APP}\) tienen,
respectivamente, firmas \((3,1,0)\), \((1,0,3)\) y \((1,1,2)\).
La corona \(\{i:K_i\geq729\}=B_K=H_e\cap P_\pi\) enlaza
el registro numérico con una cuaterna de incidencia, y la cocadena
firmada fija \(H^-_\alpha=B_K\sqcup\{5,7\}\).
El censo completo sitúa la firma \((3,1,0)\) en quince cuaternas.
La rigidización corresponde al marco ordenado
\((H^-_\alpha,H_e,H_\varphi,H_{\rm APP})\): su estabilizador
en \(G_W\) es la identidad. El resultado fija la libertad de
permutación de las doce posiciones; los valores de \(K\) conservan
su determinación por el registro integral y sus operaciones de
reconstrucción.

El levantamiento entero de la firma regional añade un resultado
complementario. Las identidades \(p+e+f=q+3c\) y
\(p+e-f=a+3h\) conservan los cocientes que acompañan a los residuos,
y el cociclo de acarreo hace compatible su composición con la
asociación de la suma. En la tabla completa de \(27\) columnas,
las ternas \((0,0,0)\) y \((1,2,0)\) tienen \(q=a=0\), mientras
sus parejas \((c,h)\) son \((0,0)\) y \((1,1)\).
Cada una de las reducciones \(\bar c,\bar h\), añadida por separado
a la evaluación de \(1,p,e,f\), aumenta el rango de cuatro a cinco.
La lectura completa \((q,a,c,h)\) conserva a su vez fibras:
\((1,2,0)\) y \((2,1,0)\) producen el mismo dato
\((0,0,1,1)\). Las identidades de transporte y las cotas del cilindro
precisan, por tanto, la información adicional de la memoria entera
y la multiplicidad de antecedentes que permanece en cada lectura.
Esta articulación conserva conjuntamente el contenido numérico y
la incidencia de la historia regional.

% END S27_FRAMING


La identidad \(C^2=5I_6\) procede de la correlación cuadrática evaluada
en las cinco posiciones finitas. Saturación y neutralidad dual dejan
una pareja de fronteras cuyo representante depende de la orientación
declarada. Los controles de clausura distinguen esa selección de las
variaciones de borde y de posición del registro. 

El resultado excepcional se obtiene siguiendo la incidencia de \(K\)
hasta el álgebra de vértices. El registro dodecafásico se construye antes
de su lectura racional y de la selección de soportes. Su configuración
firmada interviene en la bandera de Witt; el código ternario determina
el pegado de \(A_2^{12}\), y la bandera con la métrica, el origen y la
cámara declarados determina el vecino de Leech. Las extensiones de
órdenes dos y tres de su álgebra reticular realizan \(V^\natural\)
por aplicación de los teoremas de construcción correspondientes, con
sus hipótesis verificadas en el corredor reticular. Los isomorfismos
elegidos transportan vacío, vector conforme, graduación, producto y
trazas; no identifican un automorfismo de orden dos con otro de orden
tres. El Monstruo actúa en \(V^\natural\), y Moonshine determina sus
trazas graduadas. La lectura racional del registro y las cifras de las
constantes conservan otros dominios: su conexión con el álgebra es la
cadena de objetos y mapas construida, no una acción sobre sus cifras.

En la realización
excepcional, la identidad racional demuestra
\(J-(t_3+12)\in3^9q\mathbb Z[[q]]\); su primer coeficiente determina
la optimalidad del exponente uniforme. La prueba formal conserva su
alcance en todos los grados y las comprobaciones truncadas conservan
su función de contraste finito.
% END R27_FRAMING_CONCLUSION_ES

El mismo \(K\) determina una polarización geométrica y una referencia
de recuperación. En la representación real fijada, el vector no nulo
\(u_K\) determina una recta dentro de \(V_{11}\); su norma cuadrática
exacta normaliza el término de rango uno del proyector ortogonal de rango
diez, cuya ley de covariancia queda especificada.
La órbita dodecafásica completa es invertible y permite identificar
operadores a partir de sus respuestas con cotas exactas de estabilidad.
Su codificación racional conserva los doce bloques con la coordenatización fijada;
el balance bilateral recupera el estado desde dos canales definidos por isometrías,
y su extensión unitaria recupera conjuntamente estado y memoria
entrante. La invariancia de otros observables conserva la condición
de conmutación probada. La función aritmético-geométrica de \(K\)
comprende, por tanto, incidencia, dirección e identificación, además
de su evaluación escalar.

La prolongación unital de la representación de Weyl establece una
realización continua de la dimensión tracial. Los operadores de fase
y desplazamiento generan cada álgebra matricial finita, y las
inclusiones de fibra completa conservan unidad y traza normalizada.
El cierre uniformemente hiperfinito nonádico tiene una representación
tracial cuyo cierre débil es un factor hiperfinito de tipo \(II_1\);
en él, los rangos normalizados alcanzan todas las dimensiones traciales
del intervalo unidad. La continuación marcada utiliza otra inclusión
y produce otro cierre. El resultado identifica la operación de
refinamiento responsable de conservar la normalización global.

\Needspace{14\baselineskip}
La acción y el par angular determinan la forma compacta, en lugar de
recibir un radio prescrito. La renovación regional normalizada y las
secciones positivas de acción fijan la escala común de la elipse;
sus coordenadas angulares fijan el cociente de semiejes. El radio
positivo único satisface
\[
 R_\star^4=\frac{A-C^*}{A+C^*}
 =\frac{499\alpha-\eta_{\mathrm{ret}}}{501\alpha+\eta_{\mathrm{ret}}}.
\]
La pertenencia a la coordenatización elíptica y la elección de la raíz positiva
están demostradas. El producto de semiejes conserva la acción;
su intercambio transforma el radio en \(R_\star^{-1}\).
Las coordenatizaciones angulares, la orientación de los ejes y la escala de acción
mantienen así sus respectivas funciones en la forma cuadrática.

La dualidad T establecida es una equivalencia escalar y operatoria.
Una misma involución conserva la forma compacta y conjuga unitariamente
los Hamiltonianos autoadjuntos de radios recíprocos, incluidos sus
dominios y los de sus formas cuadráticas. La covariancia de la sección
residual exige transportar \(L_9\); sobre el cociente correspondiente,
la energía mínima está bien definida y conserva la dualidad.
La normalización de cuerda transporta estas identidades a los radios
dimensionales sin modificar el radio de forma generado. La colección
elíptica contiene ambas fibras recíprocas y conserva la restricción
de los dominios y del isomorfismo de pantallas. La semiconjugación
modular sobre el eje imaginario enlaza sus parámetros y sus funciones,
conservando el dominio propio de los operadores de cada realización.

Las pantallas convierten la jerarquía dimensional en aplicaciones
especificadas. Las reducciones \(12\to11\to10\) y \(26\to24\)
proceden, respectivamente, de la eliminación del modo uniforme y de la
dirección \(u_K\), y de la restricción ortogonal seguida por el cociente
isotrópico del retículo lorentziano. El morfismo conjunto identifica
las presentaciones cocientadas y entrelaza el sector compacto.
El grafo conserva la coordenada en \(V_{11}\) junto con su proyección;
cada proyección aislada mantiene su núcleo. La perforación ternaria
añade la pantalla combinatoria de longitud once con distancia cinco
y partición perfecta. Los rangos quedan determinados por estas
operaciones, no por una identificación entre espacios de dimensión
coincidente.

La prolongación hacia la teoría M tiene un soporte algebraico
determinado y una condición dinámica precisa. Los cambios de dirección
y de tipo aritmético de las trayectorias, compuestos con su dato
precedente de frontera, forman un cociclo exactamente aditivo.
La superficie de mundo, la acción de Polyakov y el codominio de campos
undecadimensionales reciben después esa estructura mediante los mapas
de realización. El teorema de transporte conserva \(Q^2=H\) en la
imagen de la isometría bajo los entrelazamientos y dominios declarados.
La identificación física completa incluye los campos, las tensiones
y la normalización de acción, el sector de osciladores y sus amplitudes,
el control de anomalías y el límite gravitatorio de baja energía de
cada compactificación
(sección~\ref{iv:condiciones-realizacion-fisica}).
Estos requisitos fijan la dinámica en la que se transportan las
estructuras demostradas.

Queda establecida una cadena común que enlaza la construcción discreta
del continuo con el álgebra Moonshine, la polarización dodecafásica,
la realización tracial, la dualidad compacta y la composición de acción.
Sus mapas determinan qué información permanece, qué cociente se toma
y qué ecuaciones se conservan. Éste es el contenido matemático que
reciben las realizaciones dimensionales y sobre el que continúan los
desarrollos de estadística, radiación y estructura de las especies.

\par\endgroup
